\documentclass[11pt]{article}
\usepackage[T1]{fontenc}
\usepackage[utf8]{inputenc}
\usepackage[letterpaper,hmargin=2.4cm,vmargin=1.5cm]{geometry}
\usepackage{amsmath,amssymb}
\pagestyle{plain}
\setlength{\parskip}{4pt}
\setlength{\emergencystretch}{3em}

\begin{document}

\begin{center}
{\Large\bfseries Retraction of Theorem 4.4 and Proposition 3.8, and Status Note (v2)}\\[5pt]
{\large Almost Reflection Positivity for Gradient-Flow Observables via\\
Gaussian Localization in Lattice Yang--Mills Theory}\\[6pt]
Lluis Eriksson \quad\textbullet\quad ai.viXra:2602.0084 \quad\textbullet\quad July 2026
\end{center}

\noindent\rule{\textwidth}{0.8pt}

\vspace{2pt}
\noindent
\textbf{The principal results of this paper are withdrawn.} They rest on three
independent defects (A--C in \S2--\S4), two of which are outright false lemmas rather
than gaps. A subsequent audit of the statements this note had initially listed as
\emph{surviving} produced four more (D--G in \S6), three of them further false
statements, none of them arising from the Wilson-flow gauge-mode obstruction that
causes A--C. One of the
false lemmas had already been refuted, with an explicit counterexample, in a companion
paper of this same programme --- in July 2026 --- and the correction was never
propagated here.
\textbf{Nothing below shows that the conclusion of Theorem 4.4 is false}; it shows
that it is unproved. Section 7 records the programme's three existing repair routes,
which address defects A--C; \S6 separately records the further reformulations that
defects D--G would require, and those are not in hand. The v1 manuscript follows
unmodified.

\section{Status}

\begin{center}
\begin{tabular}{|l|l|}
\hline
\textbf{Statement} & \textbf{Status as of v2}\\
\hline
Definitions 2.2, 2.5, 2.6 & Stand, as definitions\\
\hline
Theorem 4.1 (lattice RP, Osterwalder--Seiler) & Stands (cited, not ours)\\
\hline
Lemma 3.1 (heat-kernel domination) & \textbf{FALSE}, twice over (\S3, \S6.4)\\
\hline
Lemma A.1 (variance--oscillation bound) & \textbf{FALSE} (\S4)\\
\hline
Theorem 5.1 (OS reconstruction, conditional) & \textbf{FALSE as stated} (\S6.1)\\
\hline
Lemma 3.3 (graph distance vs.\ time sep.)
 & \textbf{FALSE} on the periodic torus;\\
 & recoverable with a two-boundary\\
 & buffer or a periodic distance (\S6.2)\\
\hline
Lemma 3.5 (chain rule for oscillation)
 & \textbf{FALSE as printed}; recoverable\\
 & by integrating the Jacobian along\\
 & the single-link geodesic (\S6.3)\\
\hline
Lemma 3.6 (Gaussian tail on the link graph)
 & Recoverable only after replacing its\\
 & heat-kernel input and fixing the\\
 & volume/time regime (\S6.4)\\
\hline
Lemma 3.4 (trans-plane influence decay) & \textbf{Withdrawn}\\
\hline
Proposition 3.8 (half-space approximation) & \textbf{Withdrawn}\\
\hline
Theorem 4.4 (quantitative RP for flowed obs.) & \textbf{Withdrawn}; and see \S6.2 on\\
 & its \emph{statement}, not only its proof\\
\hline
Proposition 5.3 (OS axioms for $\omega^t_L$) & \textbf{Withdrawn} (\S2)\\
\hline
\end{tabular}
\end{center}

\section{Defect A: the object of the paper is not constructed}

Section 2 states: ``\emph{The limit state $\omega^t_L$ on $\mathcal A$ exists by
[1, Theorem 1.1]}'', and the paper's opening sentence describes $\omega^t_L$ as ``the
continuum limit state constructed in [1]''. Reference [1] is
\textbf{ai.viXra:2602.0085}, whose Theorem 1.1 has been \textbf{withdrawn} in its own
v2: the linearization identity underlying its ultraviolet estimate is false, because
the unfixed Wilson flow has undamped gauge directions.

Consequently $\omega^t_L$ is not constructed by the cited result, and no alternative
construction is provided in this paper. Every statement of this
paper \emph{about} $\omega^t_L$ --- Theorem 4.4 and Proposition 5.3 in particular ---
is a statement about an object whose construction has been retracted. Step 4 of the
proof of Theorem 4.4 invokes [1, Theorem 1.1] directly, for the convergence
$\omega^t_k((\Theta F)^*F)\to\omega^t_L((\Theta F)^*F)$.

\section{Defect B: Lemma 3.1 is the withdrawn lemma, imported}

Lemma 3.1 is stated here as ``\emph{Heat kernel domination --- Paper 83, Lemma 3.6}'',
namely
\[
\|X_\tau(\ell,\ell_0)\|_{\mathrm{op}}\le p_\tau(\ell,\ell_0),
\qquad X_\tau(\ell,\ell_0)=\partial_{U_{\ell_0}}W_\tau(U)_\ell ,
\]
a bound on the columns of the Jacobian of the Wilson flow map. This is exactly the
statement refuted in the v2 of ai.viXra:2602.0085, and the refutation applies verbatim
here. At $U\equiv\mathbf 1$ every plaquette is trivial, so the configuration is a
global minimum of the Wilson action and a fixed point of the flow. For
$\lambda:\Lambda^0_k\to\mathfrak{su}(N)$ supported at a single vertex, the gauge
variation $X^\lambda(e)=\lambda(x)-\lambda(y)$, $e=(x,y)$, is supported on the $2d$
incident links and is \emph{exactly stationary} under the flow, since the whole gauge
orbit consists of minima. By linearity it is a signed sum of the $2d$ single-link
variations $X_\tau(\cdot,e_i)$. Fix any one incident link $e_*$; the gauge mode has
unit norm there for every $\tau$, so Lemma 3.1 would force
\[
1=\bigl\|X^\lambda_\tau(e_*)\bigr\|\;\le\;\sum_{i=1}^{2d}p_\tau(e_*,e_i)
\;\xrightarrow[\ \tau\to\infty\ ]{}\;\frac{2d}{|\mathcal G_k|},
\]
the limit because the scalar heat kernel on a finite connected graph converges to its
uniform stationary value. This is impossible as soon as $|\mathcal G_k|>2d$. (Both
sides here are pointwise, as Lemma 3.1 is; the $\ell^2$ form of the same contradiction
is the one used against Proposition 3.9 in the v2 of ai.viXra:2602.0085, whose bound
is an $\ell^2$ column bound.) Lemma 3.4 is proved from Lemma 3.1 together with Lemma
3.6, and is therefore withdrawn.

\section{Defect C: Lemma A.1 is false, and this programme had already said so}

Appendix A asserts, under the heading ``\emph{Variance--oscillation bound (no product
assumption)}'', that for \emph{any} probability measure $\mu$ on $G^E$ and any
$S\subset E$,
\[
\operatorname{Var}_\mu\!\bigl(\Phi\mid\mathcal G_S\bigr)\;\le\;
\tfrac14\sum_{e\in E\setminus S}\operatorname{osc}_e(\Phi)^2 ,
\]
with the proof step ``\emph{each $\Delta_i$ is the conditional expectation (given
$\mathcal F_{i-1}$) of the difference obtained by varying $U_{e_i}$ only, so
$\|\Delta_i\|_{L^\infty}\le\frac12\operatorname{osc}_{e_i}(\Phi)$}'', and the closing
remark ``\emph{no independence or product structure on $\mu$ is needed}''.

\paragraph{The proof step is false.}
The martingale increment
$\Delta_i=\mathbb E[\Phi\mid\mathcal F_i]-\mathbb E[\Phi\mid\mathcal F_{i-1}]$ is
\emph{not} obtained by varying $U_{e_i}$ alone: conditioning on $U_{e_i}$ also changes
the conditional law of the remaining variables $U_{e_{i+1}},\dots,U_{e_n}$. Fix
$g_0\neq g_1$ in $G$ and a continuous $f:G\to[-1,1]$ with $f(g_0)=1$,
$f(g_1)=-1$ (the range restriction is what makes the oscillations below exactly $1$).
On $E=\{e_1,e_2\}$ let
$\mu=\frac12\delta_{(g_0,g_0)}+\frac12\delta_{(g_1,g_1)}$ and
$\Phi(U_1,U_2)=f(U_2)$. Then $\operatorname{osc}_{e_1}(\Phi)=0$, while
$\Delta_1=\mathbb E[\Phi\mid U_{e_1}]-\mathbb E[\Phi]=f(U_{e_1})$ has modulus $1$.

\paragraph{And so is the inequality itself.}
The example above refutes the proof step but not the conclusion: there
$\operatorname{osc}_{e_2}(\Phi)=2$ pays the whole variance, and the asserted bound
holds with equality. A direct counterexample to the \emph{inequality} is obtained from
the same measure by spreading the observable over both coordinates. With $\mu$ as
above and
\[
\Phi(U_1,U_2)=\tfrac12\bigl(f(U_1)+f(U_2)\bigr),
\]
$\Phi$ takes the values $\pm1$ with probability $\tfrac12$ each under $\mu$, so
(taking $S=\emptyset$, for which $\mathcal G_S$ is trivial)
$\operatorname{Var}_\mu(\Phi)=1$. But varying either coordinate alone moves $\Phi$ by
exactly $1$, so $\operatorname{osc}_{e_1}(\Phi)=\operatorname{osc}_{e_2}(\Phi)=1$ and
the asserted bound reads
\[
1=\operatorname{Var}_\mu(\Phi)\;\le\;\tfrac14\bigl(1^2+1^2\bigr)=\tfrac12 ,
\]
which is false. Lemma A.1 is therefore false as stated, not merely unproved.

This is not a new observation. It is the two-spin counterexample recorded in the v2 of
\textbf{ai.viXra:2602.0070}, and it is the reason the v2 of
\textbf{ai.viXra:2602.0077} withdrew its own Lemma 1.5 --- which asserted the same
oscillation control ``for any probability measure'' by the same martingale-increment
argument. Those corrections were made in July 2026. \textbf{The present paper was
submitted in February 2026 and never replaced, so the correction never reached it.}

The measure to which Lemma A.1 is applied here is $\mu_k$, the Wilson measure at scale
$k$ (Remark A.2). It is not a product measure --- that is the entire content of an
interacting lattice gauge theory. Proposition 3.8 derives its $L^2$ error bound from
Lemma A.1 and is therefore withdrawn on this ground alone, independently of \S3.

\section{What falls, and what does not}

Theorem 4.4 fails through three independent routes: its Step 1 rests on Proposition
3.8, which rests on both defective lemmas; its Step 4 rests on the withdrawn Theorem
1.1 of [1]; and its subject $\omega^t_L$ is not constructed. Proposition 5.3 is
withdrawn for the same reason as its subject.

Exactly two things are untouched:
\begin{itemize}
\item \textbf{Definitions 2.2, 2.5, 2.6}, as definitions.
\item \textbf{Theorem 4.1}: exact reflection positivity of the Wilson measure at the
lattice level, cited from Osterwalder--Seiler. Not a result of this paper, and not
affected by anything here.
\end{itemize}

\noindent
Everything else that a first reading of this note listed as surviving did not survive
the audit of \S6.

\noindent
\textbf{The conclusion of Theorem 4.4 is not refuted.} Almost-reflection-positivity
for flowed observables, with a defect exponentially small in the buffer distance, may
well hold. Nothing above bears on its truth.

\section{Defects D--G: the survivors, audited}

An earlier draft of this note carried a list of results ``untouched'' by defects
A--C. Each had been cleared by checking what it \emph{depends on}. None had had its
proof read. Reading them produced four more defects, three of them false statements.
None arises from the Wilson-flow gauge-mode obstruction of \S3: D is an error about
semigroups, E about the periodic geometry of the torus, F about differentiating a
nonlinear map, G about the long-time behaviour of a heat kernel on a finite graph.
(F does concern the flow map, but its error is generic to nonlinear maps and owes
nothing to the special structure of the Wilson flow.) They are recorded in full,
because a retraction that quietly shortens its own survivors list between drafts is
worth nothing.

\subsection{D. Theorem 5.1 is false as stated}

The theorem assumes reflection positivity on $\mathcal E_+$ and a strongly continuous
semigroup $\{\alpha_\tau\}$ of $*$-endomorphisms with (i)
$\alpha_\tau(\mathcal E_+)\subset\mathcal E_+$, (ii) $\omega\circ\alpha_\tau=\omega$,
(iii) $\Theta\circ\alpha_\tau=\alpha_\tau\circ\Theta$; and concludes, in item 3, a
\emph{self-adjoint} $H\ge 0$ with $T_\tau=e^{-\tau H}$. Step 5 of the proof deduces
this from Hille--Yosida.

Hille--Yosida gives no such thing. It produces a generator $-H$ with $H$ accretive; it
does not make $H$ self-adjoint. Self-adjointness of the OS Hamiltonian comes from the
\emph{reflection} relation between $\Theta$ and the evolution, typically
$\Theta\alpha_\tau=\alpha_{-\tau}\Theta$ in a two-sided setting, which yields
$\langle F,T_\tau G\rangle_{\mathrm{OS}}=\langle T_\tau F,G\rangle_{\mathrm{OS}}$.
Hypothesis (iii) is \emph{commutation}, a different relation, and it does not give
symmetry.

What (i)--(iii) do give is stronger than a contraction and useless for the conclusion:
since $\alpha_\tau$ is a $*$-endomorphism, (iii) and (ii) give
\[
\|T_\tau[F]\|^2_{\mathrm{OS}}
=\omega\bigl((\Theta\alpha_\tau F)^*\alpha_\tau F\bigr)
=\omega\bigl(\alpha_\tau((\Theta F)^*F)\bigr)
=\omega\bigl((\Theta F)^*F\bigr)=\|[F]\|^2_{\mathrm{OS}},
\]
so $T_\tau$ is an \emph{isometry}. A semigroup of isometries fixing $\Omega_{\rm vac}$
has no decay to extract.

\paragraph{Counterexample.} Take $\mathcal A=\mathcal E_+=M_2(\mathbb C)$,
$\Theta=\mathrm{id}$, $\omega=\frac12\operatorname{tr}$, and
$\alpha_\tau(F)=U_\tau^*FU_\tau$ with $U_\tau=e^{i\tau K}$ for a self-adjoint
non-scalar $K$. Reflection positivity holds ($\omega(F^*F)\ge0$), and (i)--(iii) hold.
The OS space is $M_2(\mathbb C)$ with the Hilbert--Schmidt inner product, and
$T_\tau=e^{-i\tau\,\mathrm{ad}_K}$ is a \emph{unitary group}. It equals $e^{-\tau H}$
with $H\ge0$ self-adjoint only if $H=0$, i.e.\ only if $K$ is scalar. Conclusion 3 of
Theorem 5.1 fails. Separability (conclusion 1) also does not follow for an arbitrary
$*$-algebra.

\paragraph{What survives of it.} The final sentence of the theorem --- that applying
the construction to the lattice Wilson measure $\mu_k$ yields a lattice Hilbert space
and $H_k\ge0$ unconditionally --- is true, but \emph{not by this theorem}. It is the
Osterwalder--Seiler transfer-matrix result already quoted here as Theorem 4.1, where
self-adjointness and positivity of the transfer operator are established by the
reflection structure, not by Hille--Yosida. Theorem 5.1 is recoverable by assuming
directly that the induced $T_\tau$ are self-adjoint contractions, or by replacing
(iii) with the reflection relation above and proving symmetry from it. Either way a
new proof is required, and the statement as printed must not be used.

\subsection{E. Lemma 3.3 is false on the periodic torus}

Remark 2.1 fixes time coordinates on $\mathbb T^4_L$ by the representative
$x_0\in(-L/2,L/2]$, and $x_0(\ell)$ is the minimum of the representatives of the two
endpoints of $\ell$. Lemma 3.3 asserts that adjacent links satisfy
$|x_0(\ell)-x_0(\ell')|\le a_k$, its proof reasoning that ``the time coordinates of
the endpoints differ by at most $a_k$''. Across the periodic cut they do not.

Let $v$ be a lattice site whose time representative $x_*$ is \emph{maximal} in
$(-L/2,L/2]$. (One should not assume $L/2$ is itself a site: if the number of temporal
sites is odd, that representative does not occur. Maximality is all that is needed,
and it holds for either parity.) The spatial link $\ell_s$ based at $v$ has both
endpoints at time $x_*$, so $x_0(\ell_s)=x_*$. The temporal link $\ell_t$ leaving $v$
in the positive time direction wraps across the cut, since $x_*+a_k>L/2$ by
maximality, so its other endpoint has representative $x_*+a_k-L$ and
$x_0(\ell_t)=x_*+a_k-L$. The two links share $v$, hence are adjacent, yet
\[
|x_0(\ell_s)-x_0(\ell_t)|=L-a_k .
\]
The consequence drawn from the lemma,
$\operatorname{dist}_{\mathrm{gr}}(\ell,\ell_0)\ge|x_0(\ell)-x_0(\ell_0)|/a_k$, fails
with it: these two links are at graph distance $1$.

The underlying issue is not a slip in a two-line proof. \textbf{On a torus the
positive-time region has two boundaries}, the plane $x_0=0$ and the cut
$x_0=\pm L/2$; the support condition $x_0>\delta$ separates an observable from the
first and not from the second. Recovery requires a two-sided buffer,
$\delta<x_0(\ell)<L/2-\delta'$, or a formulation in terms of periodic distance and
distance to the complement of the positive region. \textbf{This bears on the
\emph{statement} of Theorem 4.4, not only on its proof}: its support hypothesis
$\{x_0>\delta_0\}$ is, as written, insufficient on the torus.

The manuscript is aware of this phenomenon elsewhere --- Remark 5.2 observes that on
$\mathbb T^4_L$ ``literal time translation by large $\tau$ can wrap around and need
not preserve the half-space algebra'' --- and does not apply the same care in
Lemma 3.3.

\subsection{F. Lemma 3.5 is false as printed}

The proof of Lemma 3.5 bounds, for $U,U'$ differing only at $\ell_0$,
\[
d_k\bigl(W_{\tau_k}(U),W_{\tau_k}(U')\bigr)\le
\frac{1}{|\Lambda^1_k|^{1/2}}
\Bigl(\sum_{\ell}\|X_{\tau_k}(\ell,\ell_0)\|_{\mathrm{op}}^2\,
d_G(U_{\ell_0},U'_{\ell_0})^2\Bigr)^{1/2},
\]
where $X_{\tau_k}(\ell,\ell_0)$ is, by Definition 2.3, the differential of the flow
\emph{at a configuration}. This is the estimate valid for a linear map, applied to a
nonlinear one. For a nonlinear map the mean-value inequality requires the derivative
along the connecting path, not at one endpoint: writing $U^{(s)}$ for the geodesic in
the $\ell_0$ variable from $U_{\ell_0}$ to $U'_{\ell_0}$, of length $\rho$, the
correct bound carries
$\int_0^{\rho}\|X_{\tau_k}(\ell,\ell_0)(U^{(s)})\|_{\mathrm{op}}\,ds$, or the cruder
$\rho\,\sup_V\|X_{\tau_k}(\ell,\ell_0)(V)\|_{\mathrm{op}}$. Neither the integral nor
the supremum appears; the supremum taken at the end of the proof is over the endpoint
$U_{\ell_0}$ only.

In fairness this defect alone is repairable, and would have been invisible in the
original architecture: the (false) Lemma 3.1 asserted a bound on
$\|X_\tau(\ell,\ell_0)\|_{\mathrm{op}}$ uniform in the configuration, so inserting the
supremum would have cost nothing. It is recorded because the lemma as printed does not
hold, and because the repair is no longer available once Lemma 3.1 is gone.

\subsection{G. The scalar bound in Lemma 3.1 omits the stationary term}

Independently of the gauge modes of \S3, the right-hand side of Lemma 3.1 is itself
false for large $\tau$. It asserts, for the heat kernel of a finite graph and
\emph{all} $\tau>0$,
\[
p_\tau(\ell,\ell_0)\le\frac{C}{(\tau+1)^{d/2}}
\exp\Bigl(-\frac{c\,\operatorname{dist}_{\mathrm{gr}}(\ell,\ell_0)^2}{\tau+1}\Bigr).
\]
On a finite connected graph $p_\tau(\ell,\ell_0)\to1/|\mathcal G_k|>0$ as
$\tau\to\infty$, while the right-hand side tends to $0$. The bound therefore fails at
large times before any question about the Wilson flow arises.

The companion paper ai.viXra:2602.0085 carries the corresponding bound \emph{with} the
stationary term $1/|\mathcal G_k|$, and devotes its Remark 3.10 to the size of that
term. The present paper dropped it. Consequently Lemma 3.6, whose proof consumes this
bound shell by shell, cannot be recovered merely by renaming its kernel: recovery
requires restoring the stationary term, or restricting to times before mixing, or
using the correct periodic kernel with its images --- and, for localization against
one half of the torus, handling the two boundaries of \S6.2.

\section{Three repairs, all of which exist in this programme}

\begin{enumerate}
\item[(a)] \textbf{For Lemma A.1.} The programme's own repair is the Dobrushin-type
decoupling hypothesis (H-DEC/AT) of ai.viXra:2602.0070 and 2602.0072 v2, under which
the increment is controlled by the oscillations with an influence-leakage constant.
ai.viXra:2602.0077 v2 already adopted it. Applying it here converts Proposition 3.8
into a conditional statement with an explicit constant.
\item[(b)] \textbf{For Lemma 3.1.} The routes recorded in ai.viXra:2602.0085 v2: prove
the oscillation bound for the damped flow $\alpha_0>0$ (for gauge-invariant $F$ the
composed observable is unchanged, so the bound transfers), or estimate transverse to
the gauge orbit on the quotient.
\item[(c)] \textbf{For the subject.} Repair of Theorem 1.1 of ai.viXra:2602.0085, on
which the existence of $\omega^t_L$ depends.
\end{enumerate}

\noindent
Each needs a new proof. Note that (a) and (b) are independent: fixing either one alone
leaves Proposition 3.8 unproved.

\section{On how this happened}

The mathematical defects in \S3 and \S4 are ordinary errors. What is worth recording
is \S4's history: the claim was refuted inside this programme, in a companion paper,
and the refutation did not propagate to a paper that depended on it. The papers were
corrected; the dependency graph between them was not maintained. That is the same
failure that produced the archive's metadata errata --- no manifest, no propagation
step --- and it is being addressed by a replacement ledger that records, for each
record, what depends on it.

A second lesson is recorded because it was learned the hard way in the drafting of
this very note. Theorem 5.1 appeared in the first draft's table as a survivor. It had
been cleared by checking what it \emph{depends on}, not by reading its proof; the
defect in \S6 sits in that proof and has nothing to do with the gradient flow. In a
retraction, the list of survivors deserves exactly the adversarial scrutiny given to
the list of casualties --- a survivor is a positive claim being made under the
author's name at the moment when the reader's trust is lowest, and it is the one place
where the temptation to salvage is strongest.

\vspace{4pt}
\noindent\rule{\textwidth}{0.8pt}

\noindent\footnotesize
\textbf{Provenance.} Defect B follows from an external adversarial review of
ai.viXra:2602.0085, verified independently against that manuscript. Defects A and C
were found on 2026-07-29 by tracing the dependencies of that retraction through the
companion papers, exactly the propagation step whose absence is described in \S8. A
further round of the same external review supplied Defect D, corrected the norm
mismatch in the displayed contradiction of \S3, and pointed out that the original
counterexample in \S4 refuted the proof step of Lemma A.1 without refuting its
conclusion; the direct counterexample to the inequality was added in response, and
each point was verified against the manuscript before being adopted. A later
survivor audit, in the same review, supplied defects E--G: the periodic-cut
counterexample to Lemma 3.3, the missing path integration in Lemma 3.5, and the
omitted stationary term in the finite-graph heat-kernel bound; a subsequent pass
replaced the witness of \S6.2 by the maximal-representative site, so that it no longer
presumes $L/2$ to be a lattice point. Each was checked directly against the unchanged
v1 proof before inclusion. It is recorded that three of the four survivor-audit
defects came from outside. No
compiled theorem in the Lean core of \texttt{THE-ERIKSSON-PROGRAMME} discharges any
hypothesis of this paper, and none is affected by this retraction. Remark 4.3 of the
v1 manuscript records that the Hamiltonian used in the companion papers is
reconstructed from exact lattice reflection positivity (Theorem 4.1), not from the
almost-RP estimate withdrawn here; that statement is unaffected, and limits how far
this retraction propagates.

\end{document}
